\documentclass[12pt]{article}

% Lei de atracao universal de Newton
% Gambelas, Tue Oct 15 14:32:35 WEST 2013
% Written by Tordar

\usepackage{graphicx}
\usepackage[svgnames]{xcolor}
\usepackage{pst-all}
\usepackage{pst-3dplot}
\usepackage{pst-slpe}
\usepackage{pifont}



\begin{document}

\pagestyle{empty}

\begin{pspicture}(-1,-1)(9,5)
\pscircle[linewidth=0.1mm,
linecolor=white,
fillstyle=ccslope,
slopebegin=white,
slopeend=DarkGreen,
slopecenter=0.65 0.65,
slopesteps=25](0,0){1.5}
\pscircle[linewidth=0.1mm,
linecolor=white,
fillstyle=ccslope,
slopebegin=white,
slopeend=DimGray,
slopecenter=0.65 0.65,
slopesteps=25](8,0){0.75}
\psline[linewidth=0.8mm,linecolor=black,linestyle=solid]{<->}(0,-2)(8,-2)
\psline[linewidth=1mm,linecolor=blue,linestyle=solid]{->}(0,0)(2,0)
\psline[linewidth=1mm,linecolor=blue,linestyle=solid]{<-}(6,0)(8,0)
\rput{0}(4,-1.7){\scalebox{1.4}{$r$}}
\rput{0}(0,1.7){\scalebox{1.4}{$m_1$}}
\rput{0}(8,1.0){\scalebox{1.4}{$m_2$}}
\end{pspicture}

\end{document}
